\documentclass[11pt]{article}
\usepackage[margin=1in]{geometry}
\usepackage{amsmath,amssymb,amsthm}
\usepackage[T1]{fontenc}
\usepackage{lmodern}
\usepackage[hidelinks]{hyperref}
\hypersetup{pdftitle={Droop-core nonemptiness for approval committees with two omissions},pdfauthor={Jiarui Fang},pdfsubject={Approval-based committee elections and coalitional stability},pdfkeywords={approval voting, Droop core, Hare core, PAV, core existence}}
\newtheorem{theorem}{Theorem}
\newtheorem{lemma}[theorem]{Lemma}
\newtheorem{corollary}[theorem]{Corollary}
\newcommand{\PAV}{\operatorname{PAV}}
\newcommand{\gain}{\operatorname{Q}}
\title{Droop-core nonemptiness for approval committees\\with two omissions}
\author{Jiarui Fang\\[0.35em]
\small Boston University\\
\small \href{mailto:baymin@bu.edu}{\texttt{baymin@bu.edu}}\\
\small ORCID: \href{https://orcid.org/0009-0006-9100-0445}{0009-0006-9100-0445}}
\date{September 2026}
\begin{document}
\maketitle
\begin{abstract}
Core stability requires that no coalition can use its proportional share
of seats to make every member strictly better off. We prove that every
finite approval election with $m\ge3$ candidates and $k=m-2$ seats has a
deterministic Droop-core committee, and hence a Hare-core committee,
without restricting the number of voters or their approval sets.
The Droop rule can let smaller coalitions fund a challenge, making
stability more demanding than under Hare.
Individual improvements are strict, and equality at the Droop threshold
does not block. The proof bounds repeated
support for blocking targets by the size of their intersection. This
estimate justifies a contraction that fixes an initial blocking target,
removes its strict gainers and leaves a smaller ordinary election with
two omissions. Every Droop core of the residual election lifts to a
Droop core of the original election. Induction gives a finite exact
construction, with overlapping deviations and an explicit zero-seat
boundary. Together with a one-swap bound, this establishes Droop-core
existence whenever at most two candidates are unselected.
\end{abstract}

\section{Introduction}
A representative committee can be challenged by voters who would all
prefer to use their collective entitlement to select a different set.
The Hare core rules out every such challenge when a coalition's share
of seats is proportional to its share of the electorate. With binary
approvals, each voter's utility is the number of approved candidates
selected. The resulting stability requirement is stronger than a
comparison of total approval scores: each member of a blocking
coalition must strictly improve.

Proportional Approval Voting (PAV) assigns diminishing harmonic credit
for each additional approved representative. Peters' one-swap bound gives
core existence when $k=m$ or $k=m-1$~\cite[Lemma~3.1(ii)]{Peters2025}.
He also proved existence for $k\le8$ and, using a recursive PAV
procedure, for $m\le15$~\cite[Theorems 4.1, 4.5 and 5.3]{Peters2025}.
His procedure fixes blocking targets and removes the voters who strictly
gain from them before
reoptimizing over full committees~\cite[Algorithm 1]{Peters2025}.
Becker, Greger and Peters establish existence with at most seven distinct
voter types by a computer-assisted argument~\cite[Theorems 2.5 and 5.2]{BGP2026}.
The result here extends the omission-based guarantee to $m-k=2$ under
the stronger Droop criterion, with no bound on candidate count or voter types.
There is no restriction on how many candidates a voter approves or
disapproves. This quota distinction matters: Peters shows that PAV can
fail the Droop core already at $k=6$~\cite[Example~6.1]{Peters2025}.

The main step is a contraction theorem rather than a claim that PAV
itself always selects a core committee. Start from a global PAV
maximizer and a blocking target $F$. If some of its strict gainers can
strictly gain again after $F$ is selected, two omissions force their
approval sets into a specific configuration. A summed swap inequality
then bounds their number in terms of overlap with the later target.
This bound permits a new ordinary residual election on $C\setminus F$,
with fewer voters and $k-|F|$ seats. Every residual core lifts, so no
additional invariant or tie-breaking condition is required.

The proof is analytic. We state the Droop result first, then develop
the contraction through a direct Hare proof before proving the stronger statement.
Exact examples illustrate the distinction between the quotas, the contraction and its
repeated-support bound; the computations are not premises of the proof.

\section{Model and main theorem}\label{sec:model}
Let $N$ be a nonempty finite voter set, $n=|N|$, and let $C$ be a finite
candidate set with $m=|C|$. Voter $i$ approves an arbitrary set $A_i\subseteq C$.
For every $X\subseteq C$, put $u_i(X)=|A_i\cap X|$. A committee has exactly
$k$ candidates, where $1\le k\le m$. A nonempty coalition $S\subseteq N$
Hare-blocks a committee $W$ by a target $T\subseteq C$ if
\[
 n|T|\le k|S|\quad\text{and}\quad
 u_i(T)>u_i(W)\quad\text{for every }i\in S.
\]
An empty target cannot strictly improve a voter, and affordability implies
$|T|\le k$. Define the full strict-gainer set and its size by
\[
 \gain(W,T)=\{i\in N:u_i(T)>u_i(W)\},\qquad q(W,T)=|\gain(W,T)|.
\]
Every blocking coalition is contained in $\gain(W,T)$. Conversely, if a
nonempty target satisfies $kq(W,T)\ge n|T|$, its full strict-gainer set is
nonempty and is itself an affordable coalition. Thus the exact core test is
\begin{equation}\label{eq:core}
 W\text{ is Hare core}\quad\Longleftrightarrow\quad
 kq(W,T)<n|T|\quad
 \text{for every }\varnothing\ne T\subseteq C,\ |T|\le k.
\end{equation}
Equality is blocking under this Hare convention. Under the Droop
convention of Peters~\cite[Section~6]{Peters2025}, affordability instead
requires $n|T|<(k+1)|S|$. Taking the full strict-gainer set again shows
that a nonempty target Droop-blocks $W$ exactly when
\begin{equation}\label{eq:droop}
 (k+1)q(W,T)>n|T|.
\end{equation}
A Droop-core committee admits no such target. Individual improvement
remains strict, but equality in \eqref{eq:droop} does not block.
Since $q(W,T)\le n$, an affordable integral target has size at most $k$.
Every Droop-core committee is Hare core: for a nonempty target,
$kq\ge n|T|$ implies $q>0$ and hence $(k+1)q>n|T|$.
Neither convention requires unanimous approval of a target, and targets
may overlap the committee. Unqualified uses of \emph{core} and
\emph{blocker} in Sections~\ref{sec:repeat}--\ref{sec:construction}
refer to Hare, unless both conventions are explicitly mentioned.

\begin{theorem}\label{thm:main}
For every finite approval election with $m\ge3$ candidates and $k=m-2$ seats,
a deterministic Droop-core committee exists. In particular, the
ordinary deterministic Hare core is nonempty.
\end{theorem}
The Droop statement is proved in Section~\ref{sec:droop}.

Proportional Approval Voting (PAV) maximizes a harmonic approval score.
We use its optimality only to initialize a contraction; a PAV maximizer
need not itself be core. Define
\[
 H_0=0,\qquad H_a=\sum_{j=1}^a\frac1j\quad(a\ge1),\qquad
 \PAV(X)=\sum_{i\in N}H_{u_i(X)}.
\]
A global PAV maximizer among $k$-committees exists by finiteness. It is also
a one-swap local maximizer: replacing any selected candidate by any
unselected candidate does not increase PAV.

\section{Bounding repeated support}\label{sec:repeat}
We first record an elementary removal bound. Fix a committee $W$, an
unselected candidate $b$, and a nonempty set $Y\subseteq W$. For one voter,
sum the PAV changes over the swaps $W-\{y\}+\{b\}$, $y\in Y$. If $b$ is
approved, every such change is nonnegative. If $b$ is not approved and
$u=u_i(W)>0$, the sum is $-|A_i\cap Y|/u\ge-1$. If $b$ is not approved
and $u=0$, the sum is zero.
Thus this sum is at least $-1$ for every voter; no division by zero is used.

The following numerical bound implies the one-outsider exclusion of
Peters~\cite[Lemma~3.1(ii)]{Peters2025} under both quota conventions.

\begin{lemma}[One-outsider bound]\label{lem:one}
Let $W$ be a one-swap PAV local maximizer with $|W|=k$.
For every nonempty $T\subseteq C$ with $|T|\le k$ and
$|T\setminus W|\le1$,
\begin{equation}\label{eq:one}
 (k+1)q(W,T)\le n|T|.
\end{equation}
Thus such a target is neither a Hare blocker nor a Droop blocker.
If exactly two candidates are unselected, every blocker under either
convention contains both of them.
\end{lemma}
\begin{proof}
If $T\subseteq W$, then $q(W,T)=0$, so \eqref{eq:one} holds.
Otherwise write $T\setminus W=\{b\}$.
Write $t=|T|$ and $Y=W\setminus T$, so $d=|Y|=k-t+1\ge1$. A strict gainer
approves $b$ and none of $Y$. Its utility at $W$ is at most $t-1$; hence
each of the $d$ swaps adding $b$ and removing a member of $Y$ increases its
PAV term by at least $1/t$. Every other voter's summed contribution is at
least $-1$ by the removal bound. Writing $q=q(W,T)$, the total change,
summed over all $d$ swaps and voters, is therefore at least
\[
 \frac{dq}{t}-(n-q)=\frac{(k+1)q-nt}{t}.
\]
Local maximality makes each swap's total change nonpositive, so
$(k+1)q\le nt$. This excludes Droop blocking directly.
If $q>0$, it also gives $kq<(k+1)q\le nt$; if $q=0$, then
$kq=0<nt$. Thus Hare blocking is also excluded, proving the consequences.
\end{proof}

The next estimate holds for \emph{every} completion of a target containing
both initially unselected candidates. Neither affordability of that target
nor optimality of the later committee is needed.

\begin{lemma}[Repeated-support bound]\label{lem:repeat}
Let $|C|=k+2$, let $W_0$ be a one-swap PAV local maximizer, and let
$C\setminus W_0\subseteq F\subseteq C$ with $|F|\le k$.
Put $Q_0=\gain(W_0,F)$. For any $k$-committee
$W\supseteq F$ and any $T\subseteq C$ with $|T|\le k$, define
\[
 I=F\cap T,\quad s=|I|,\quad
 R=Q_0\cap\gain(W,T),\quad r=|R|.
\]
Then
\begin{equation}\label{eq:repeat}
 (k+1)r\le ns.
\end{equation}
The later target $T$ is not required to be affordable.
\end{lemma}
\begin{proof}
If $r=0$, the conclusion is immediate. Suppose $r>0$ and put
$B=C\setminus W_0$, so $|B|=2$. For $i\in R$, let $a=|A_i|$. The two
strict improvements and the containment $F\subseteq W$ give
\[
 a\ge u_i(T)\ge u_i(W)+1\ge u_i(F)+1
 \ge u_i(W_0)+2\ge a.
\]
Every inequality is therefore an equality. In particular, the voter
approves both members of $B$, $F$ contains exactly $a-1$ of its approvals,
and $T$ contains all its approvals. By assumption $B\subseteq F$,
so $B\subseteq I$. The unique approved candidate $d_i$ outside $F$ lies
in $W_0\setminus I$, and every other approved candidate lies in $I$.
Consequently $a\ge3$ and $a-1\le s$, in particular $s\ge2$.

Set $Y=W_0\setminus I$. Since both unselected candidates belong to $I$,
\[
 d=|Y|=k-s+2.
\]
Sum the PAV changes over all $2d$ swaps adding $b\in B$ and removing
$y\in Y$. For a voter $i\in R$, both additions are approved, exactly one
member of $Y$ is approved, and $u_i(W_0)=a-2$. This voter's summed
contribution is exactly
\[
 \frac{2(d-1)}{a-1}\ge\frac{2(k-s+1)}s.
\]
For each voter outside $R$, the removal bound applied separately to the
two choices of $b$ gives a summed contribution of at least $-2$. Every
swap has nonpositive aggregate change, so
\[
 0\ge\frac{2r(k-s+1)}s-2(n-r).
\]
Multiplication by the positive number $s/2$ and rearrangement yield
$(k+1)r\le ns$, as required.
\end{proof}

\section{A core-preserving contraction}\label{sec:contraction}
\begin{lemma}[Contraction and lifting]\label{lem:lift}
Let $|C|=k+2$, let $W_0$ be a global PAV maximizer, and let $F$ block $W_0$.
Write
\[
 f=|F|,\quad Q_0=\gain(W_0,F),\quad q_0=|Q_0|,\qquad
 D=C\setminus F,\quad N'=N\setminus Q_0,
\]
\[
 k'=k-f,\qquad n'=n-q_0.
\]
Then $k'\ge1$, $n'\ge1$, $|D|=k'+2$, and $kn'\le nk'$. Consider the
ordinary residual approval election with electorate $N'$, candidates $D$,
approval sets $A_i\cap D$, and $k'$ seats. Every core committee $U$ of
this residual election lifts to an original core committee $F\cup U$.
\end{lemma}
\begin{proof}
First, $q_0<n$: otherwise any padding of $F$ to $k$ candidates would
strictly improve every voter over $W_0$, and hence strictly increase PAV,
contrary to global maximality. The original blocking inequality
$kq_0\ge nf$ then also implies $f<k$. Thus $n',k'>0$ and
\begin{equation}\label{eq:residualquota}
 kn'=k(n-q_0)\le n(k-f)=nk'.
\end{equation}
The candidate count satisfies $|D|=k'+2$.

Let $U$ be a residual core committee and set $W=F\cup U$. Suppose a
nonempty target $T$ of size $t\le k$ blocks $W$ in the \emph{original}
election, so $kq\ge nt$ for $q=q(W,T)$. Define
\[
 I=F\cap T,\quad s=|I|,\quad V=T\setminus F,\quad v=|V|,
\]
\[
 r=|Q_0\cap\gain(W,T)|,\qquad q'=q-r.
\]
Lemma~\ref{lem:one} gives $C\setminus W_0\subseteq F$, so
Lemma~\ref{lem:repeat} gives $(k+1)r\le ns$, hence $kr\le ns$.
It follows that
\begin{equation}\label{eq:remaining}
 kq'=kq-kr\ge n(t-s)=nv.
\end{equation}
The target $V$ is nonempty, since $T\subseteq F\subseteq W$ would make
strict improvement impossible. Every voter $i\in\gain(W,T)\setminus Q_0$
strictly improves $V$ over $U$ in the residual election: cancellation of
the fixed part gives
\[
 |A_i\cap V|>|A_i\cap U|+|A_i\cap(F\setminus T)|
 \ge |A_i\cap U|.
\]
Combining \eqref{eq:residualquota} and \eqref{eq:remaining},
\[
 kk'q'\ge nk'v\ge kn'v,
 \qquad\text{so}\qquad k'q'\ge n'v.
\]
As $q'\le n'$ and $n'>0$, this inequality implies $v\le k'$.
Thus $V$ is a nonempty admissible target and the $q'$ remaining voters
form an affordable strictly improving residual coalition, contradicting
the core property of $U$. Any additional residual strict gainers only
strengthen this contradiction.
\end{proof}

The residual election is a separately defined induction instance; its
quota is not substituted into an original blocking test. Every original
test retains the original $n,k$. Equations~\eqref{eq:residualquota}--
\eqref{eq:remaining} prove exactly why a hypothetical original blocker
would induce a residual blocker.

This contraction differs from fixed-target PAV reoptimization. With
$F$, $N'$ and $U\subseteq D$ as above, the two objectives are
\[
 \underbrace{\sum_{i\in N'}H_{|A_i\cap F|+|A_i\cap U|}}
 _{\text{full-committee score with }F\text{ fixed}}
 \qquad\text{and}\qquad
 \underbrace{\sum_{i\in N'}H_{|A_i\cap U|}}
 _{\text{ordinary residual-election score}}.
\]
Peters' recursive PAV rule uses the first objective~\cite[Algorithm~1]{Peters2025};
our recursive initialization uses the second. They do not generally differ
by a constant independent of $U$. More importantly, Lemma~\ref{lem:lift}
accepts every residual core, irrespective of how it is obtained.

\section{Induction and exact construction}\label{sec:construction}
We first prove the Hare-core consequence of Theorem~\ref{thm:main}
directly; Section~\ref{sec:droop} establishes its stronger Droop statement.
\begin{proof}[Direct proof of Hare-core nonemptiness]
Induct on $k\ge1$, always with $|C|=k+2$. For $k=1$, choose a global
PAV maximizer. Any blocker would have size one and $q=n$, hence strictly
increase PAV, a contradiction.

For $k>1$, choose a global PAV maximizer $W_0$. If it is core, the claim
is proved. Otherwise choose any blocking target $F$. Lemma~\ref{lem:lift}
produces a residual ordinary approval election with a nonempty electorate,
$k'+2$ candidates and $1\le k'<k$ seats. The induction hypothesis gives
a residual core committee $U$. Lemma~\ref{lem:lift} lifts it to the
original core committee $F\cup U$, which has exactly $f+k'=k$ candidates.
\end{proof}

The proof gives the following exact construction.
\begin{center}
\fbox{\begin{minipage}{0.92\linewidth}
\textbf{HareCore}$(N,C,k,(A_i)_{i\in N})$, where $|C|=k+2$.
\begin{enumerate}
\item Choose a global PAV maximizer $W_0$ among the $k$-committees.
\item If no target blocks $W_0$ under \eqref{eq:core}, return $W_0$.
Otherwise choose any blocker $F$.
\item Set $N'=N\setminus\gain(W_0,F)$, $D=C\setminus F$ and $k'=k-|F|$.
\item Compute $U=\textbf{HareCore}(N',D,k',(A_i\cap D)_{i\in N'})$
and return $F\cup U$.
\end{enumerate}
\end{minipage}}
\end{center}
The positive integer seat count strictly decreases at each recursive call.
Any tied PAV maximizer and any first blocker are allowed. At a call with
$m$ candidates, PAV maximization compares $\binom{m}{2}$ committees;
the exhaustive blocking test examines
$\sum_{t=1}^{m-2}\binom mt=2^m-m-2$ targets.
Thus the reference implementation is not a polynomial-time algorithm.

All lemmas sum over actual voters, so repeated types with positive integer
multiplicities are included. Empty approvals make zero PAV changes and
never strictly gain; universal approvals cannot strictly gain from a target
of at most $k$ candidates. These voters are nevertheless retained in every
electorate unless they genuinely belong to a deleted full-gainer set.
Finiteness is used for global maximization, exhaustive tests and the
well-founded induction.


\subsection{Contraction of a blocked PAV maximizer}\label{sec:example}
Let $C=B\mathbin{\dot\cup}X\mathbin{\dot\cup}Y$, where
$B=\{b_0,b_1\}$, $|X|=3$ and $|Y|=6$. There are $18$ voters:
five approve $X\cup\{b_0\}$, five approve $X\cup\{b_1\}$, and eight
approve $Y$. Set $k=9$.

The committee $W_0=X\cup Y$ is the unique global PAV maximizer.
To see this, measure the PAV loss of a committee relative to the full
candidate set $C$. Every committee omits a pair, and its loss depends
only on the pair's classes:
\[
\begin{array}{c|rrrrrr}
\text{omitted-pair classes}&BB&BX&BY&XX&XY&YY\\ \hline
\PAV(C)-\PAV(W)&\frac52&\frac{25}6&\frac{31}{12}&
\frac{35}6&\frac{23}6&\frac{44}{15}
\end{array}
\]
The $BB$ loss is strictly smallest. Nevertheless,
$F=X\cup B$ blocks $W_0$: the first ten voters improve from utility
three to four, and
\[
 kq(W_0,F)=9\cdot10=18\cdot5=n|F|.
\]
For $k=9$, this example has fewer voters than the $27$-voter
unique-global-PAV counterexample described as voter-minimal in
Peters~\cite[Section 4.3]{Peters2025}. We make no minimality claim for
eighteen voters. This comparison is specific to $k=9$: a PAV failure with
$k=18$ and three voters appears in Becker, Greger and
Peters~\cite[Figure~1]{BGP2026}.

After contraction, eight voters remain, all approving the six candidates
in $Y$, and the residual committee size is four. Every four-element
subset $U\subseteq Y$ is residual core because no target of at most
four candidates can strictly improve any voter. Lemma~\ref{lem:lift}
therefore makes $F\cup U$ an original core committee.
The accompanying code reproduces this construction and independently
checks every original target, using integers for all quota tests.
Here the deleted voters approve no candidate outside $F$, so repeated
support is zero. For a nonzero example, keep $B,X,Y,k$ but change the
three multiplicities to $27,27,44$ and add one voter approving
$B\cup\{y_1\}$, where $y_1\in Y$.
The resulting $99$-voter profile again has the unique PAV maximizer
$W_0=X\cup Y$ and the Hare blocker $F=B\cup X$, now with
$9\cdot55=99\cdot5$.
Choose any four-element $U\subseteq Y$ excluding $y_1$.
For the added voter,
\[
 u_i(W_0)=1,\qquad u_i(F)=u_i(F\cup U)=2,\qquad
 u_i(B\cup\{y_1\})=3.
\]
Thus a deleted voter can strictly gain again. The later target
$T=B\cup\{y_1\}$ has $q(F\cup U,T)=1$ and is not affordable under
either quota: $9\cdot1<99\cdot3$ and $10\cdot1<99\cdot3$.
The example illustrates the scope of Lemma~\ref{lem:repeat}, not a
blocker of the lifted committee. Its validity does not require deleted
voters to be permanently unable to gain. The accompanying checker
replays the bound for all completions of $F$ and all legal targets.

\section{The Droop strengthening}\label{sec:droop}
We now use the strict Droop test \eqref{eq:droop} throughout the induction.
\begin{proof}[Proof of Theorem~\ref{thm:main}]
For induction only, also allow $k=0$ and $m=2$, with a nonempty
electorate. The empty committee is then Droop core: no nonempty
target can satisfy \eqref{eq:droop}. This is an auxiliary base case,
not a change to the positive-seat model of Section~\ref{sec:model}.

Let $k\ge1$ and choose a global PAV maximizer $W_0$. If it is Droop
core, stop. Otherwise choose a Droop blocker $F$, and put
$f=|F|$, $Q_0=\gain(W_0,F)$ and $q_0=|Q_0|$.
Lemma~\ref{lem:one} implies that $F$ contains both candidates outside $W_0$,
and Lemma~\ref{lem:repeat} applies. Moreover $q_0<n$, since otherwise
padding $F$ to $k$ candidates strictly increases every voter's PAV term.

Define $D=C\setminus F$, $N'=N\setminus Q_0$, $k'=k-f$ and
$n'=n-q_0$. These satisfy $n'>0$, $0\le k'<k$, $|D|=k'+2$, and
\begin{equation}\label{eq:drooptransfer}
 (k+1)n'<n(k'+1).
\end{equation}
By induction the ordinary residual election on $N',D$ with projected
approvals and $k'$ seats has a Droop core $U$. We show that every such
$U$ lifts to a Droop core $W=F\cup U$.

Suppose instead that $T$ Droop-blocks $W$. Write
$q=q(W,T)$, $t=|T|$, $s=|F\cap T|$, $V=T\setminus F$, $v=|V|$,
$r=|Q_0\cap\gain(W,T)|$ and $q'=q-r$. Since $T\subseteq F$
cannot strictly improve anyone over $W$, we have $v\ge1$.
Lemma~\ref{lem:repeat} and \eqref{eq:droop} give
\[
 (k+1)q'=(k+1)(q-r)>n(t-s)=nv.
\]
Using \eqref{eq:drooptransfer} and only positive integer multipliers,
\[
 (k+1)(k'+1)q'>n(k'+1)v>(k+1)n'v,
 \qquad\text{hence}\qquad (k'+1)q'>n'v.
\]
The $q'$ voters outside $Q_0$ strictly improve from $U$ to $V$ after
cancelling the fixed part, exactly as in Lemma~\ref{lem:lift}.
Also $q'\le n'$ and $n'>0$, so $v<k'+1$. Integrality gives $v\le k'$.
For $k'=0$ this contradicts $v\ge1$; for $k'>0$ it makes $V$ an
admissible Droop blocker of $U$, again a contradiction.
The lift has exactly $f+k'=k$ candidates, completing the induction.
\end{proof}

The same exhaustive construction therefore works with the Droop test
\eqref{eq:droop} and an empty-committee return at the auxiliary zero-seat
base case. The original Hare construction and its equality convention
remain as stated in Section~\ref{sec:construction}.

\begin{corollary}\label{cor:atmosttwo}
Every finite approval election with $1\le k\le m$ and $m-k\le2$
has a deterministic Droop-core committee, and hence a Hare-core committee.
\end{corollary}
\begin{proof}
For $m-k\le1$, choose a global PAV maximizer $W$.
Every legal target has $|T\setminus W|\le1$, so Lemma~\ref{lem:one}
excludes Droop blocking. The case $m-k=2$ is Theorem~\ref{thm:main}.
\end{proof}

\subsection{A Hare core outside the Droop core}\label{sec:distinction}
Peters' Example~6.1~\cite{Peters2025} separates the two quota conventions
within the two-omission class. Relabel its candidates as
$C=B\mathbin{\dot\cup}X\mathbin{\dot\cup}Y$, with
$B=\{b_0,b_1\}$, $|X|=2$, $|Y|=4$ and $k=6$.
The $24$ voters have the following approvals:
\[
\begin{array}{c|r}
\text{approval set}&\text{multiplicity}\\ \hline
X\cup\{b_0\}&7\\
X\cup\{b_1\}&7\\
Y&10
\end{array}
\]
The unique global PAV maximizer is $W_0=X\cup Y$.
For omitted-pair classes $BB$, $BX$, $BY$, $XX$, $XY$ and $YY$, the respective losses
relative to $\PAV(C)$ are $14/3,49/6,29/6,35/3,43/6,35/6$;
the $BB$ loss is uniquely smallest.
For $F=X\cup B$, both seven-voter types strictly improve, and
\[
 6\cdot14=84<24\cdot4=96<7\cdot14=98.
\]
Thus $F$ Droop-blocks $W_0$ but does not Hare-block it.
In fact, $W_0$ is Hare core. The ten voters approving $Y$ are
satiated. A target improving exactly one seven-voter type needs at
least three candidates; improving both requires all four candidates of $F$.
The Hare tests fail in both cases, since $6\cdot7<24\cdot3$ and
$6\cdot14<24\cdot4$.

Consequently the Hare construction returns $W_0$, whereas the Droop
construction fixes $F$ and leaves ten voters, four candidates and two seats.
Every pair $U\subseteq Y$ is a residual Droop core and lifts to
$F\cup U$. The supplied input and independent verifier reproduce the
two outputs and check all $246$ original legal targets for each.

\subsection{Why Droop equality must remain nonblocking}
The strict inequality in \eqref{eq:droop} cannot be replaced by a
weak inequality in the existence theorem. For any $k\ge1$, take
$n=k+1$ voters and $C=\{c_1,\ldots,c_{k+1},d\}$.
Voter $i$ approves only $c_i$, and no voter approves $d$.
Every $k$-committee $W$ omits some $c_i$. For $T=\{c_i\}$,
only voter $i$ strictly gains, so
\[
 (k+1)q(W,T)=k+1=n|T|.
\]
Hence the modified convention that blocks at Droop equality has an
empty core even with two omissions. This is a boundary of the quota
convention, not an empty-core example under \eqref{eq:droop} or a
tightness claim for the repeated-support bound.

\section{Discussion}
The two-omission condition is used at one precise point. A voter who
strictly improves first to $F$ and then to a target against a completion
of $F$ must start exactly two approvals below its total and gain one
approval at each step. This rigidity turns repeated participation into
the overlap bound of Lemma~\ref{lem:repeat}. Counting two possible gains
per voter alone would not justify deleting a coalition. The overlap
bound supplies the missing charge for candidates cancelled from a
later target.

With three or more omissions, two strict improvements no longer force
the utility equalities used in Lemma~\ref{lem:repeat}. The proof therefore
does not establish core existence for unrestricted $(m,k)$. Its conclusion
is a uniform integral existence theorem for the entire two-omission class,
under both Hare and Droop affordability, supported by contraction and lifting.

\section*{AI use disclosure}

OpenAI Codex, using GPT-6 Astra, was used for proof exploration and
checking, exact-verifier development, and manuscript drafting and revision.
AI-generated output is not treated as mathematical evidence. The results
are established by the analytic proofs presented above. The accompanying
code implements the construction and independently checks its outputs and
the illustrative examples. The author remains
responsible for the statements, proofs, citations, source code, and final
artifacts.

\section*{Code availability}

The constructor and verifier source code, example instances, manuscript
source, and reproduction instructions accompany this manuscript. The
project repository is
\url{https://github.com/Baymax-ray/approval-core-two-omissions}.

\begin{thebibliography}{2}
\bibitem{Peters2025}
D. Peters, \emph{The Core of Approval-Based Committee Elections with Few Seats},
Proceedings of the Thirty-Fourth International Joint Conference on
Artificial Intelligence (IJCAI 2025), pp.~4014--4022.
\href{https://doi.org/10.24963/ijcai.2025/447}{doi:10.24963/ijcai.2025/447}.
Full version: arXiv:2501.18304v2, 14 May 2025.
\url{https://arxiv.org/abs/2501.18304v2}.
\bibitem{BGP2026}
P. Becker, M. Greger and D. Peters,
\emph{Core Existence in Approval-Based Committee Elections with up to Seven
Voter Types}, arXiv:2605.06194v2, 17 August 2026.
\url{https://arxiv.org/abs/2605.06194v2}.
\end{thebibliography}
\end{document}
